% !TEX program = latexmk
% !TEX encoding = UTF-8 Unicode

% 注意：请使用 TeXLive 2020 及以上版本，已在 2020 2021 版本测试编译成功，已在 2019 测试编译失败
% 如不想更新，可尝试使用 手动编译 XeLaTeX（可能出错）
% WARNING: Please use TeXLive version >= 2020, This template has been tested and found available on 2020, 2021, unavailable on 2019. 

% 注意：请确保自己已经完整阅读了 README.md 这一 markdown 文档，部分功能的使用方法并未在 .tex 中直接给出；如果仍有使用问题，请在 Github 上提出 issue
% WARNING: Pleans read README.md first cause some usages are not given in .tex files.

%%%%%%%%%%%%%%%%%%%%%%%%%%%%%%%%%%%%%%%%%%%%%%%%%%%%%%%%%%%%%%%%%%%%%%%%%%%%%%%%%%%%%%%%%%%%%%%%%%%%%
%% select the basic style of this thesis/dissertation (this document will use thesis for convenience)
%% 选择论文的基本类型

\documentclass[
    doctor,     % 必选项：   {master, doctor} 此处不区分专业/学术学位，在下面学位类型处区分
                % Mandatory: {master, doctor} No difference between Academic degree and Professional degree,
                %                             but be careful about arguments in `\degree' and `\subject'
    % english,  % 可选项：   英文正文请选择此项
                % Optional:  For english main content, It will change some auto-generated matter into English
    % blind,    % 可选项：   论文用于盲审请选择此项，将不会产生致谢(内容)、答辩委员会会议决议(内容)、
                %            常规评阅人名单(内容)，注意请自行修改 `攻读学位期间取得的研究成果' 内的内容格式
                % Optional:  For blind review, It will not generate Acknowledgements(Content), Decision of Defense Committee(Content),
                %            General Reviewers List(Content), CHANGE THE STYLE of Achievements BY YOURSELF
]{XJTU-thesis}



%%%%%%%%%%%%%%%%%%%%%%%%%%%%%%%%%%%%%%%%%%%%%%%%%
%% fill the each blank for auto-generate contents
%% 填写以下信息用以自动生成

% 论文标题（不超过35个字，英文注意大小写规律）
% Title, Make sure you have an acceptable capitalization
\title{这是一个不超过三十五个字的名字比较长的关于如何开展XXX的研究}{English Title and English Title and English Title}

% 学位类型
% 考虑到专业学位的学位名称的少数特例(如专业型法学硕士不是 Master of Juris 而是 Juris Master)，此处学位类型请按照案例和文件填写
% 同时为了前向兼容，可选项「A/P」代表「学术/专业」型学位，默认为学术型
% Type of your degree, Translate it from documents in `Materials/Requirements/2021/01 中英文题名页示例/英文标准翻译/'
% \degree[A]{硕士}{Master of Engineering} % 学术型(Academic)硕士请基于 '学术学位名称.txt' 填写
% \degree[P]{工程硕士}{Master of Engineering} % 专业型(Professional)硕士请基于 '专业学位（领域）英文标准翻译.pdf' 填写
%\degree[A]{博士}{Doctor of Philosophy}  % 学术型(Academic Doc)博士请填写 '{博士}{Philosophy}'
\degree[P]{工程博士}{Doctor of Engineering} % 专业型(Professional Doc)博士请基于 '专业学位（领域）英文标准翻译.pdf' 填写

% 作者姓名（注意：所有人名英文均为「名在前，姓在后」，如果只有外文名，请在两个参数都填写外文名称）
% If you have only foreign name, put it as both first and second argument
\author{郑正正}{Zhengzheng Zheng}

% 指导老师姓名（注意：同作者姓名）
% Name of supervisor, It have the same requirments as the author
\advisor{冯峰}{教授}{Feng Feng}{Prof.}

% 合作指导老师姓名 或 老师团队（合作指导老师指：1.与招生简章中一致的合作导师，2.CSC项目的合作导师）
%（校方模板要求只能选择一个，都有则显示合作导师）
% Name of associate advisor or adviror's team, You can use only one of them, and advisorassociate has a higher priority.
% \advisorassociate{陈尘}{副教授}{Chen Chen}{Asscociate Prof.}
% \advisorteam{团队中文名}{English Name of the Team}

% 学科名称，请基于 '学科（专业）英文标准翻译.pdf' 填写
% Name of the subject, also get it from that file
% \subject{航空宇航科学与技术}{Aeronautical and Astronautical Science and Technology}
\subject{航空工程}{Aeronautical Engineering}

% 答辩委员会委员 显示的顺序和这里的一样 第一个人是主席
% Committee member of your defence, notice that the order shows in the thesis is same as here, and the first one is the chairperson
% each member should be put as {Organization,Name,Title} split by comma
\addcommitteemember{西安交通大学,张长长,教授}
\addcommitteemember{西安理工大学,王旺旺,教授}
\addcommitteemember{国网陕西经济技术研究院,李力,高工}
\addcommitteemember{西安交通大学,东方不败,副教授}
\addcommitteemember{西安交通大学,赵照,研究员}

% 答辩时间（手动指定）
% Defence date, input manually
\defensedate{2021}{06}{22}

% 答辩地点（涉密论文请手动设置为「西安交通大学」）
% Defence location, default value is 「西安交通大学」
\defenseloc{西安交通大学主楼E座303室}

% 论文提交日期，不输入参数则默认使用当前日期，如手动指明年月，请在第一个可选参数内填写年份，第二个可选参数填写月份（均为阿拉伯数字）
% Submission date of this thesis, if you not put it manually, it will use the current time
% \submitdate[2021][06]
\submitdate

% 常规审阅人 要求和答辩委员会委员一样
% General reviewer list, same requirements as the committee member
\addgeneralreviewer{西安交通大学,张长长,教授}
\addgeneralreviewer{西安理工大学,王旺旺,教授}
\addgeneralreviewer{国网陕西经济技术研究院,李力,高工}
\addgeneralreviewer{西安交通大学,东方不败,副教授}
\addgeneralreviewer{西安交通大学,赵照,研究员}

% 学院
% School or Faculty, unused now
% \school{电气工程学院}{School of Electrical Engineering}

% 专业[学士学位使用]
% Major, unused now
% \stumajor{计算机科学与计数}{Computer Science}

% 学号[学士学位使用]
% Student ID, unused now
% \stuid{}

% 班级[学士学位使用]
% Administrative class, unused now
% \adminclass{电气7xx班}{}

% 参考文献源 参数中不要添加 .bib
% 请使用 \addreferenceresource 添加数据库(可导入多个参考文献数据库)
% 若自动化导入攻读学位期间的成果则使用 \addachivementresource
\addreferenceresource{References/reference}
\addreferenceresource{References/reference}
\addachivementresource{References/achievement}

%%%%%%%%%%%%%%%%%%%%%%%%%%%%%%%%%%%%%%%%%%%%
%% 如果有使用其他包，请在这里添加
%% If you need other packages, use them here

% \usepackage{}


%%%%%%%%%%%%%%%%%%%%%%%%%%%%%%%%%%%%%%%%%%%%%%%%%%%%%%%%%%%%%%%%%%%%%%%%%%%%%%%%%%%
%% 注意：根据校方要求，以下所有页面顺序不可调整
%% Notice: The order of these pages are defined in the requirements of the University.

%% 但在最终提交前，可以通过注释所有 \thesis 开始的命令设置是否生成各个部分
%% 或根据说明调整 latexmkrc 文件，使用 \includeonly 命令导入部分章节。

% \includeonly{
%   Main_Spine/c1,
%   Main_Spine/c2,
%   Main_Spine/c3,
%   Main_Spine/c4,
%   Main_Spine/c5,
%   Main_Spine/c6,
% }

\begin{document}
% [自动生成] 中英题名页
% [Auto Generate] Chinese English Title Page
\thesistitles

% [自动生成] 答辩委员会页
% [Auto Generate] Defense Committee Pages
\thesiscommittes

% 生成摘要页 修改 Main_Miscellaneous/abstract_chs/eng.tex 中的内容
% Abstract, Rewrite your content in Main_Miscellaneous/abstract_chs/eng.tex
\thesisabstract

% [自动生成] 中英目录页
% [Auto Generate] Table of Contents
\thesistableofcontens

% 主要符号表 修改 Main_Miscellaneous/glossary.tex 中的内容
% Glossaries Page, Rewrite your content in Main_Miscellaneous/glossary.tex
\thesisglossarylist

% 正文 注意英文正文写作时，中、英标题还是先中后英标题；同时，下面参数的顺序有意义，不要乱放
% Main contents, X of cX is the chapter of the thesis, you can change it if you want, the order matters
% and KEEP Chinese Title as the FIRST argument, English Titile as the SECOND, if you use `english' option

% 可以通过 \thesisbody 直接导入各部分正文，但也可通过 \thesisbodybegin & \include & \thesisbodyend 组合导入正文

% \thesisbody{
%     Main_Spine/c1,
%     Main_Spine/c2,
%     Main_Spine/c3,
%     Main_Spine/c4,
%     Main_Spine/c5,
%     Main_Spine/c6
% }

\thesisbodybegin
% !TeX root = ../main.tex

\xchapter{绪论}{Introductions}

\xsection{背景}{Backgrounds}

本文档仅提供部分可能用到的示例，系统的学习 \LaTeX 用法请参考其他书目，如 \emph{lshort}，可通过在 shell(命令行/cmd) 中执行 \clist{texdoc lshort-zh} 获得中文版。

同时，务必意识到 \LaTeX 的工作及使用方式与 Microsoft Word 有极大的不同，不要问「怎么在 \LaTeX 中实现 Word 的 xx 功能」。

\xsection{符号说明}{Examples}

为了更加清楚的说明本模板的用法，特意使用了如下几种标记方式来标记「\LaTeX 源码」及「说明」：

\clist{行内的灰色字块}表示简短的代码（\LaTeX 或 shell）；

\clist[columbiablue]{行内的蓝色字块}表示宏包的名称。

\begin{tcolorbox}
  灰色框用来放置有一定长度的 \LaTeX 源码（及执行结果）
\end{tcolorbox}

\begin{tcolorbox}[colback=red!5!white,colframe=red!75!black]
  红色框用来放置学校的学位论文模板的\textbf{格式}要求
\end{tcolorbox}

\begin{tcolorbox}[colback=blue!5!white,colframe=blue!75!black]
  蓝色框用来放置学校的学位论文模板的\textbf{内容}要求
\end{tcolorbox}

\begin{tcolorbox}[colback=blue!5!white,colframe=blue!75!black,title=绪论部分的要求]
  绪论部分主要论述论文的选题意义及应用背景、国内外研究现状分析及论文的主要研究内容等。
\end{tcolorbox}

\clearpage

\xsection{基本功能}{Basic Functions}

\xsubsection{文档层级}{Level of this document}

按照要求，正文最多应具有七个层级。其中，为了方便中英目录的生成，模板重新定义了前三级标题的命令，并新增了第六级的命令。注意，不论正文主体语言是否为英文，标题顺序均为先中后英（前三级）。

\begin{table}[H]
  \begin{tblr}{
    colspec={X[c]X[c]},
    verb
  }
    \toprule
      层级 & 命令 \\
    \midrule
      0  & \verb|\xchapter{Chs}{Eng}| \\
      1  & \verb|\xsection{Chs}{Eng}| \\
      2  & \verb|\xsubsection{Chs}{Eng}| \\
      3  & \verb|\subsubsection{Chs/Eng}| \\
      4  & \verb|\paragraph{Chs/Eng}| \\
      5  & \verb|\subparagraph{Chs/Eng}| \\
      6  & \verb|\subsubparagraph{Chs/Eng}| \\
    \bottomrule
  \end{tblr}
\end{table}

示例如下：

\subsubsection{第四级标题$1$}

\paragraph{第五级标题$1$}

\paragraph{第五级标题$2$}

\paragraph{第五级标题$3$}

\subparagraph{第六级标题$1$}

\subparagraph{第六级标题$2$}

\subparagraph{第六级标题$3$}

第六级下的一些文字

\subsubparagraph{第七级标题$1$}

第七级下的一些文字

\clearpage

\xsubsection{加粗，斜体与其他字体设置}{Bold, Emph and Other Font Settings}

在中文学位论文的写作中，不推荐使用加粗，斜体等方式突出重点。

使用\clist{\textbf{text}}进行可以对字加粗；同时，请注意到在英文中，正确的强调方式不是使用加粗，而是使用斜体，即使用\clist{\emph{text}}进行强调。如：
\begin{texcode}[]{}
  为了突出\textbf{重点}，英文可以这样：This is the \emph{key question}.
\end{texcode}

使用\clist{\textit{text}}、\clist{\underline{text}}等实现斜体，下划线等格式。





\xsubsection{脚注及其使用}{Footnotes}

\begin{tcolorbox}[colback=red!5!white,colframe=red!75!black]
  脚注可用小号字（一般小五号宋体）列在相应正文同一页最下部并与正文部分用细线（版面宽度的1/4长）隔开。
\end{tcolorbox}

脚注\footnote{脚注序号“\ding{172}，……，\ding{180}”的字体是“正文”，不是“上标”，序号与脚注内容文字之间空$1$个半角字符，脚注的段落格式为：单倍行距，段前空$0$磅，段后空$0$磅，悬挂缩进$1.5$字符；中文用宋体，字号为小五号，英文和数字用Times New Roman字体，字号为$9$磅；中英文混排时，所有标点符号（例如逗号“，”、括号“（）”等）一律使用中文输入状态下的标点符号，但小数点采用英文状态下的样式“.”。}\footnote{脚注序号“\ding{172}，……，\ding{180}”的字体是“正文”，不是“上标”，序号与脚注内容文字之间空$1$个半角字符，脚注的段落格式为：单倍行距，段前空$0$磅，段后空$0$磅，悬挂缩进$1.5$字符；中文用宋体，字号为小五号，英文和数字用Times New Roman字体，字号为$9$磅；中英文混排时，所有标点符号（例如逗号“，”、括号“（）”等）一律使用中文输入状态下的标点符号，但小数点采用英文状态下的样式“.”。}\footnote{脚注序号“\ding{172}，……，\ding{180}”的字体是“正文”，不是“上标”，序号与脚注内容文字之间空$1$个半角字符，脚注的段落格式为：单倍行距，段前空$0$磅，段后空$0$磅，悬挂缩进$1.5$字符；中文用宋体，字号为小五号，英文和数字用Times New Roman字体，字号为$9$磅；中英文混排时，所有标点符号（例如逗号“，”、括号“（）”等）一律使用中文输入状态下的标点符号，但小数点采用英文状态下的样式“.”。}是对文中有关内容的解释、说明或补充，使用上角标（序号\ding{172}、\ding{173}、$\cdots$）进行标注。

脚注可以通过 \clist{\footnote{}}自动编号生成，也可使用\clist{\footnotetex{text}}手动添加脚注。

摘要中关于项目资助的标识，是使用脚注生成的：

\begin{tcolorbox}
  \lstinline|\footnotetext{*本研究得到某某基金（编号：）的资助}|
\end{tcolorbox}


\clearpage

\xsection{数字、公式和定理环境}{Equation and Theorem}

\xsubsection{数字与单位}{Numbers and Units}

模板引入了\clist[columbiablue]{siunitx}宏包实现数字与单位的正确排版，具体使用细节请自行查看该宏包，即 \clist{texdoc siunitx}。

为了输入不同格式及角度的数字，可以使用\clist{\num{3.45d-4}}得到\num{3.45d-4}，使用\clist{\num{-e10}}得到\num{-e10}，使用\clist{\ang{1;2;3}}得到\ang{1;2;3}，避免了手动调节格式的问题。

该宏包预先定义了部分单位，可以直接调用，同时生成较为复杂的单位。如：

使用\clist{\SI{10}{\hertz}}得到 \SI{10}{\hertz}。

\xsubsection{公式、矩阵与数学符号}{Equations, Matrix and Mathematical Symbols}


\begin{tcolorbox}[colback=red!5!white,colframe=red!75!black]
  \begin{enumerate}[leftmargin=0.5cm]
    \item 公式应另起一行，居中编排，较长的公式尽可能在等号后换行，或者在“+”、“-”等符号后换行。公式中分数线的横线，长短要分清，主要的横线应与等号取平。
    \item 公式后应注明编号，公式号应置于小括号中，如公式（2-3）。写在右边行末，中间不加虚线；
    \item 公式下面的“式中：”两字左起顶格编排，后接符号及其解释；解释顺序为先左后右，先上后下；解释与解释之间用“；”隔开。
    \item 公式中各物理量及量纲均按国际标准（SI）及国家规定的法定符号和法定计量单位标注，禁止使用已废弃的符号和计量单位。
  \end{enumerate}
\end{tcolorbox}

此部分是 \LaTeX 中比较有趣的一部分，具体的使用细节请自行查阅资料。

\begin{texcode}[]{}
  公式如下：
  \begin{equation}
      -e^{\max}_\text{dis} \leq a_t \leq e^{\max}_\text{ch}\label{equation:c1:mdl:cstr_dis}
  \end{equation}
  式中：$e$为xxx；$a$为xxx。
\end{texcode}

所以如\cref{equation:c1:mdl:cstr_dis}所示：
\begin{equation}\label{equation:c1:exp}
    -e^{\max}_\text{dis} \leq a_t \leq e^{\max}_\text{ch}
\end{equation}

其中，数学公式可以使用\clist{\boldmath..\unboldmath}对整体进行加粗：

\begin{texcode}[]{}
\boldmath
\begin{equation}
  \biggl(\int_{-\infty}^\infty e^{-x^2}\diff x\biggr)^2 
\end{equation}
\unboldmath
\end{texcode}

如果只需对某个元素调整，可以通过\clist{\boldsymbol}，或使用\clist[columbiablue]{unicode-math}提供的，以\clist{\sym}开始的命令族，如\clist{\symup, \symbfup, \symit, \symbfit}：
\begin{texcode}[]{}
  \begin{equation}
    \biggl(\int_{-\infty}^\infty \symbfit{e}^{-\symup{x}^2}\diff \symbfup{x}\biggr)^2 
  \end{equation}
\end{texcode}
不建议再使用标准 \LaTeX 及 amsmath 扩展提供的 \clist{\mathbf}、\clist{\mathcal} 等输入不同样式的字母，可能会导致使用文本字体而非数学字体(的粗体形式)。

请注意，本模板已经定义了\clist{\diff}命令作为微分符号，请不要自己使用新的符号，使用效果如下：
\begin{align}\label{equation:c1:diff}
  \biggl(\int_{-\infty}^\infty e^{-x^2}\diff x\biggr)^2 
    &= \int_{-\infty}^\infty\int_{-\infty}^\infty e^{-(x^2+y^2)}\diff x\diff y \\
    &= \int_0^{2\pi}\int_0^\infty e^{-r^2}r\diff r\diff\theta \\
    &= \int_0^{2\pi}\biggl(-{e^{-r^2}\over2}\bigg\vert_{r=0}^{r=\infty}\,\biggr)\diff\theta \\
    &= \pi
\end{align}


同时，定义了部分用于简化输入的命令：

\begin{texcode}[sidebyside]{}
  $\seq{x}{n}$ \\ % serise
  $\iprod{i}{j}$ \\ % inner product
  $\mbs{B}$  % \boldsymbol
\end{texcode}

\xsubsection{定理相关}{Theorems}\label{section:Environments for math}

本模板通过\clist[columbiablue]{thmtools}宏包定义了以下环境，其中\textbf{证明}环境默认具有缩进，没有编号，同时使用\clist{\qedhere}结束证明，其他环境样式一致。

\begin{table}[H]
  \begin{tblr}{
    colspec={*{6}{X[c]}},
    vline{4},
  }
  \toprule
    环境 & 英文名称 & 中文名称 & 环境 & 英文名称 & 中文名称 \\
  \midrule
    proof & Proof & 证明 & theorem & Theorem & 定理 \\
    axiom & Axiom & 公理 & corollary & Corollary & 推论 \\
    lemma & Lemma & 引理 & definition & Definition & 定义 \\
    example & Example & 例子 & proposition & Proposition & 命题 \\
    assumption & Assumption & 假设 & remark & Remark & 注 \\
    problem & Problem & 问题 & conjecture & Conjecture & 猜想 \\
  \bottomrule
  \end{tblr}
\end{table}


\begin{texcode}[]{}
  \begin{theorem}[勾股定理]\label{theorem:pt}
    若 $a,b$ 为直角三角形的两条直角边，$c$ 为斜边，那么 $a^2 + b^2 + c^2.$
  \end{theorem}

  \begin{proof}
    通过\ldots

    所以：
    \begin{equation*}
        G(x, y) = G(y, x).  \qedhere
    \end{equation*}
  \end{proof}
\end{texcode}

\begin{proposition}
  所以：
  \begin{equation*}
      G(x, y) = G(y, x).
  \end{equation*}
\end{proposition}

\begin{conjecture}[勾股定理]
    若 $a,b$ 为直角三角形的两条直角边，$c$ 为斜边，那么 $a^2 + b^2 + c^2.$
\end{conjecture}

\begin{axiom}[勾股定理]
    若 $a,b$ 为直角三角形的两条直角边，$c$ 为斜边，那么 $a^2 + b^2 + c^2.$ 若 $a,b$ 为直角三角形的两条直角边，$c$ 为斜边，那么 $a^2 + b^2 + c^2.$ 若 $a,b$ 为直角三角形的两条直角边，$c$ 为斜边，那么 $a^2 + b^2 + c^2.$
\end{axiom}

\begin{definition}[勾股定理]
    若 $a,b$ 为直角三角形的两条直角边，$c$ 为斜边，那么 $a^2 + b^2 + c^2.$
\end{definition}

\xsection{其他环境}{Other Environments}

\xsubsection{枚举环境}{Enumerates}

\begin{texcode}[]{}
  \begin{enumerate}
    \item 123
    \item 231421
    \item 124124
  \end{enumerate}
\end{texcode}


\xsection{正文}{Text}

\zhlipsum[1]

\xsubsection{正文}{Text}

\zhlipsum[2]

\xsubsection{正文}{Text}

\zhlipsum[3]
% !TeX root = ../main.tex

\xchapter{浮动体：图表}{Floating: Figures, Tables}

请先注意，在\LaTeX 中，图表不一定要像 Word 一样被固定在某个区域，程序会自动根据上下文为图片选择合适的位置，同时还有具有超链接的交叉引用自动定位图片位置，因此不要强求图片或表格一定处于某个位置。如非要定位，则应在浮动体环境中使用\clist{[H]}选项指定位置。

\begin{tcolorbox}[colback=red!5!white,colframe=red!75!black]
  \begin{enumerate}[leftmargin=0.5cm]
    \item 图、表、公式等一律用阿拉伯数字分章连续编号，如 图1-3、表2-1、（3-2）等。图、表、公式等与正文之间间隔0.5行。

    \item 图应有图题，表应有表题，并分别置于图号和表号之后，图号和图题应置于图下方的居中位置，表号和表题应置于表上方的居中位置。引用图或表应在图题或表题右上角标出文献来源。

    \item 若图或表中有附注，采用英文小写字母顺序编号，附注写在图或表的下方。
  \end{enumerate}
\end{tcolorbox}

\xsection{图}{Figures}


\begin{tcolorbox}[colback=red!5!white,colframe=red!75!black]
  \begin{enumerate}[leftmargin=0.5cm]
    \item 插图须紧跟文述。在正文中，一般应先见图号及图的内容后再见图，一般情况下不能提前见图，特殊情况须延后的插图不应跨节；
    \item 提供照片应大小适宜，主题明确，层次清楚，金相照片一定要有比例尺；
    \item 图应具有“自明性”，即只看图、图题和图例，不阅读正文，就可理解图意。图中的标目是说明坐标轴物理意义的项目，它是由物理量的符号或名称和相应的单位组成。物理量的符号由斜体字母标注，单位的符号使用正体字母标注，量与单位间用斜线隔开。例如：$I/\SI{}{A}$，$\rho/\SI{}{kg\cdot m^{-3}}$，$F/\SI{}{N}$，$v/\SI{}{m\cdot s^{-1}}$等等；
    \item 图中用字为五号，如排列过密，用五号字有困难时，可小于五号字，但不得小于七号字；
    \item 图尺寸的\textbf{一般宽高比}应为 \SI{6.67}{cm}$\times$\SI{5.00}{cm}。特殊情况下，也可为 \SI{9.00}{cm}$\times$\SI{6.75}{cm}，或 \SI{13.5}{cm}$\times$\SI{9.00}{cm}。总之，一篇论文中，同类图片的大小应该一致，编排美观、整齐；
    \item 一幅图如有若干幅\textbf{分图}，均应编分图号，用(a)，(b)，(c),...... 按顺序编排，且各分图的分题注直接列在各自分图的正下方，总题注列在所有分图的下方正中「
  \end{enumerate}
\end{tcolorbox}

\xsubsection{单张图片的使用}{Single Figure}

\begin{texcodeonly}[]{}
  \begin{figure}[h]
    \centering
    \includegraphics[width=.5\textwidth]{xjtu.pdf}
    \caption{校徽}
  \end{figure}
\end{texcodeonly}

\begin{figure}[h]
  \centering
  \includegraphics[width=.5\textwidth]{xjtu.pdf}
  \caption{校徽}
\end{figure}

一些图片下方的文字。一些图片下方的文字。一些图片下方的文字。一些图片下方的文字。一些图片下方的文字。一些图片下方的文字。一些图片下方的文字。一些图片下方的文字。一些图片下方的文字。一些图片下方的文字。一些图片下方的文字。一些图片下方的文字。一些图片下方的文字。

\xsubsection{多张图片：应该使用 subcaption 而不是 subfigure}{Use subcaption instead of subfigure}

\begin{texcodeonly}[]{}
  \begin{figure}[h]
    \begin{subfigure}[b]{0.49\linewidth}
        \centering
        \includegraphics[width=.5\textwidth]{xjtu.pdf}
        \subcaption{校名}
    \end{subfigure}
    \begin{subfigure}[b]{0.49\linewidth}
        \centering
        \includegraphics[width=.5\textwidth]{xjtu.pdf}
        \subcaption{校名}
        \label{subfigure:icon}
    \end{subfigure}
    \caption{校徽}
    \label{fig:icon}
  \end{figure}
\end{texcodeonly}

\begin{figure}[h]
  \begin{subfigure}[b]{0.49\linewidth}
      \centering
      \includegraphics[width=.5\textwidth]{xjtu.pdf}
      \subcaption{校名}
  \end{subfigure}
  \begin{subfigure}[b]{0.49\linewidth}
      \centering
      \includegraphics[width=.5\textwidth]{xjtu.pdf}
      \subcaption{校名}
      \label{subfigure:icon}
  \end{subfigure}
  \caption{校徽}
  \label{figure:icon}
\end{figure}

\xsubsection{多张子图分页}{Break multiple figures}

\newcommand{\subfigg}{
  \begin{subfigure}[b]{0.49\linewidth}
    \centering
    \includegraphics[width=.5\textwidth]{xjtu.pdf}
    \subcaption{校名}
  \end{subfigure}
}

为了简化叙述，定义一个命令 \clist{\subfigg} 代表插入一个子图：
\begin{texcodeonly}[]{}
  \newcommand{\subfigg}{
    \begin{subfigure}[b]{0.49\linewidth}
      \centering
      \includegraphics[width=.5\textwidth]{xjtu.pdf}
      \subcaption{校名}
    \end{subfigure}
  }
\end{texcodeonly}

在\clist[columbiablue]{subfig}宏包中提供了\clist{\ContinuedFloat}命令，使一个浮动体环境与其内子图的编号跟随前一个子图，此时两个浮动体环境可以分别设置标题，若具有相同的标题，或后续图的标题设置为 \clist{（续）}，则可以实现子图分页的效果。同时，模板提供了 \clist{\floatcontinue} ，通过包装\clist{\ContinuedFloat}使得子图在此命令处断页，不产生图标题，并在下一页生成图题。此两种命令均可使用，按照喜欢或需求使用。
\clearpage
\subsubsection{使用 ContinuedFloat}

\begin{texcodeonly}[]{}
  \begin{figure}[h]
    \subfigg    \subfigg
    \caption{校徽}
  \end{figure}
  
  \begin{figure}[H]
    \ContinuedFloat
    \subfigg    \subfigg
    \subfigg    \subfigg
    \caption{校徽} % 或 \caption{（续）}
  \end{figure}
\end{texcodeonly}

\begin{figure}[h]
  \subfigg    \subfigg
  \caption{校徽}
\end{figure}

\begin{figure}[H]
  \ContinuedFloat
  \subfigg    \subfigg
  \subfigg    \subfigg
  \caption{校徽} % 或 \caption{（续）}
\end{figure}

\subsubsection{使用 floatcontinue}

\begin{texcodeonly}[]{}
  \begin{figure}[h]
    \subfigg    \subfigg
    \floatcontinue{tb}
    \subfigg    \subfigg
    \subfigg    \subfigg
    \caption{校徽}
  \end{figure}
\end{texcodeonly}

\begin{figure}[h]
  \subfigg    \subfigg
  \floatcontinue{tb}
  \subfigg    \subfigg
  \subfigg    \subfigg
  \caption{校徽}
\end{figure}


\clearpage

\xsection{表}{Tables}

\begin{tcolorbox}[colback=red!5!white,colframe=red!75!black]
  \begin{enumerate}[leftmargin=0.5cm]
    \item 如某个表需要转页接排，在随后的各页上应重复表的编号。编号后跟表题（可省略）和“（续）”，如表 1-1（续），续表均应重复表头和关于单位的陈述。表格的设计应紧跟文述。表的编排一般是内容和测试项目由左至右横读，数据依序竖读，应有自明性。若为大表或作为工具使用的表格，可作为附表在附录中给出，论文中的表格参数应标明量和单位的符号；
    \item 表中各物理量及量纲均按国际标准(SI) 及国家规定的法定符号和法定计量单位标注；
    \item 一律使用三线表，与文字齐宽，顶线与底线线粗为 $1.5$ pt，中线线粗是 $1$ pt；
    \item 使用他人表格须注明出处。
    \item 表中用字为五号字体。如排列过密，用五号字有困难时，可小于五号字，但不小于七号。
    \item 表格必须通栏，即表格宽度与正文版面平齐。
  \end{enumerate}
\end{tcolorbox}


\xsubsection{普通表格}{Normal Tables}


考虑到使用成本、内容与样式分离。模板默认调用现代化的表格绘制宏包 \clist[columbiablue]{tabularray}。该宏包的使用方法请使用 \clist{texdoc tabularray}，网上也有此手册的中文翻译。请确保使用的发行版包含这一宏包（即更新至 2023 及以后的版本）。

非跨页表格应使用 tblr 环境制作表格，使用 table 作为浮动体环境，使用 \clist{\toprule}，\clist{\midrule} 和 \clist{\bottomrule} 指定三线表的三条线：

\begin{texcode}[]{}
    \begin{table}[H]
        \caption{tblr 基本用法}
        \label{tbl:2-1}
        \begin{tblr}{
            % colspec={X[r,m,2]X[c,3]X[l,1]}
            X[r,m,2]X[c,3]X[l,1]
        }
            \toprule
            Head & Test & 测试 \\
            \midrule
            Class that has many accessor methods and accesses a lot of external data & ATFD is more than a few & 3 \\
            1 & 2 & 3 \\
            1 & 2 & 3 \\
            \bottomrule
        \end{tblr}
    \end{table}
\end{texcode}

tblr 环境各列的对齐方式可以以键值对的方式在必选参数中设置，或者直接作为唯一必选参数（参考 tabularray 手册第 24 页）。列类型包括 c/l/r, t/b/m, X 等。c/l/r 为默认的三种对齐方式，t/b/m 指定单元格内容超过列宽时垂直方向的对齐方式， X 类型会自动调整列宽使全表宽度与页宽一致，X 类型列的水平、垂直对齐方式以类似 X[c] 的方式在 X 后的方括号内指定。X 类型后的方括号可以指定此列宽度与其他X类列宽度的相对大小，也可使用 \clist{wd=A}直接指定列宽。推荐使用 X 作为列类型。

\begin{texcode}[]{}
    \begin{tblr}{
        colspec={|ll|X[c]|rr|},
        cell{1}{1}={r=3,c=2}{r,m},
        cell{1}{3}={r=2,c=3}{l,t},
        cell{3}{Z}={}{red},
        hline{4-6}={2pt,blue},
    }
        \hline
        1-1 & 1-2 & 1-3 & 1-4 & 1-5 \\
        \hline
        2-1 & 2-2 & 2-3 & 2-4 & 2-5 \\
        \hline
        3-1 & 3-2 & MIDDLE & 3-4 & 3-5 \\
        % \hline (4th)
        4-1 & 4-2 & 4-3 & 4-4 & 4-5 \\
        % \hline (5th)
        5-1 & 5-2 & 5-3 & 5-4 & 5-5 \\
        % \hline (6th)
    \end{tblr}
\end{texcode}

tabularray 的优点在于格式和内容分离，使用更贴合人类习惯。单元格的合并，文字、背景、边框的样式都可以放在参数中设置。 \clist{cell{R}{C}={merge}{style}}用于设置单元格合并，其中 R,C 指定合并单元格左上角的位置，merge参数指定合并多少行多少列，style 调整合并后的单元格样式，其中 R,C 除了使用行号进行索引，可以使用 X,Y,Z 索引倒数第三、二、一行/列，可以使用 odd/even 选择 奇/偶列，还可以使用连字符(dash)选择区间，非常方便。

\xsubsection{复杂表格}{Complicated tables}

在三线表中可以加辅助线，以适应较复杂表格的需要，具体使用方法请自行阅读宏包：

\begin{texcode}[]{}
\begin{table}[H]
  \centering
  \caption{compare with other approachs}
  \label{table:methodcompare}
  \begin{tblr}{
      colspec={*{7}{X[c]}},
      cell{1}{1}={r=2}{c,m},
      cell{1}{2,5}={c=3}{c,m},
      hline{2,4,6}={1pt},
      hline{1,3,Z}={1.5pt},
      row{4,5}={rowsep=0pt},
      row{Z}={m,font=\bfseries},
  }
    Model & trigger inditification & & & Event Extraction & & \\
    & P(\%) & R(\%) & F1(\%) & P(\%) & R(\%) & F1(\%) \\
    Baseline1 & 76.84 & 76.84 & 76.84 & 76.84 & 76.84 & 76.84 \\
    Baseline2  & 76.84 & 76.84 & 76.84 & 76.84 & 76.84 & 76.84 \\
    Baseline3  & 76.84 & 76.84 & 76.84 & 76.84 & 76.84 & 76.84 \\
    Our approach  & 76.84 & 76.84 & 76.84 & 76.84 & 76.84 & 76.84 \\
  \end{tblr}
\end{table}
\end{texcode}

\xsubsection{长表格}{Long Tables}

当表格过长时，使用\clist{longtblr}环境，此时不需要浮动体环境。在可选参数中需要设置\clist{theme=xjtu_std}，以调用预先设置好的续表样式，同时设置表格标题、label、脚注等，具体使用方法请自行阅读宏包。

\begin{texcodeonly}[]{}
    \begin{longtblr}[
        theme = xjtu_std,
        caption = {A Long Long Long Long Long Long Long Table },
        entry = {Short Caption},
        label = {tblr:test},
        note{a} = {脚注},
        note{$\dag$} = {还是脚注},
        remark{Note} = {注释，即表格内没有标注符号的脚注},
        remark{Source} = {还是注释},
    ]{
        colspec = {XX[c]X},
        rowhead = 1, rowfoot = 0,
    }
    \toprule
    Head & Head & Head \\
    \midrule
    Epsilon & Zeta\TblrNote{a} & Eta \\
    Iota & Kappa\TblrNote{$\dag$} & Lambda \\
    Nu & Xi & Omicron \\
    \dots 
    Foot & Foot & Foot \\
    \bottomrule
\end{texcodeonly}

\begin{longtblr}[
        theme = xjtu_std,
        caption = {A Long Long Long Long Long Long Long Table },
        entry = {Short Caption},
        label = {tblr:test},
        note{a} = {脚注},
        note{$\dag$} = {还是脚注},
        remark{Note} = {注释，即表格内没有标注符号的脚注},
        remark{Source} = {还是注释},
    ]{
        colspec = {XX[c]X},
        rowhead = 1, rowfoot = 0,
    }
    \toprule
    Head & Head & Head \\
    \midrule
    Head & Head & Head \\
    Alpha & Beta & Gamma \\
    Epsilon & Zeta\TblrNote{a} & Eta \\
    Iota & Kappa\TblrNote{$\dag$} & Lambda \\
    Nu & Xi & Omicron \\
    Rho & Sigma & Tau \\
    Phi & Chi & Psi \\
    %
    % ......
    Alpha & Beta & Gamma \\
    Epsilon & Zeta & Eta \\
    Iota & Kappa & Lambda \\
    Nu & Xi & Omicron \\
    Rho & Sigma & Tau \\
    Phi & Chi & Psi \\
    % Foot & Foot & Foot \\
    Alpha & Beta & Gamma \\
    \midrule
    Epsilon & Zeta & Eta \\
    Iota & Kappa & Lambda \\
    Nu & Xi & Omicron \\
    Rho & Sigma & Tau \\
    Phi & Chi & Psi \\
    Foot & Foot & Foot \\
    Alpha & Beta & Gamma \\
    Epsilon & Zeta & Eta \\
    Iota & Kappa & Lambda \\
    Nu & Xi & Omicron \\
    Rho & Sigma & Tau \\
    Phi & Chi & Psi \\
    Foot & Foot & Foot \\
    Alpha & Beta & Gamma \\
    Epsilon & Zeta & Eta \\
    Iota & Kappa & Lambda \\
    Nu & Xi & Omicron \\
    Rho & Sigma & Tau \\
    Phi & Chi & Psi \\
    Foot & Foot & Foot \\
    \bottomrule
\end{longtblr}

% !TeX root = ../main.tex

\xchapter{算法与代码}{Algorithm and Code}

\begin{tcolorbox}[colback=red!5!white,colframe=red!75!black]
  学校并未对 算法 或 代码 的排版格式作出限定，因此可以自行修改。
\end{tcolorbox}

\xsection{算法}{Algorithm}

在仔细研究后，本模板使用了基于 \clist[columbiablue]{algorithm2e} 宏包的算法排版工具，而没有使用 \clist[columbiablue]{algorithmicx (algpseudocode)}。同时，需要提醒的是， \clist[columbiablue]{algpseudocode} 在制作竖直定界线方面有一些问题，而能够较好的实现算法自动分段；\clist[columbiablue]{algorithm2e} 则相反，且这两个宏包相互冲突。

算法的具体编写方式请自行查阅宏包，即\clist{texdoc algorithm2e}，本模板已预先定义了Do-While循环即\clist{\Do{cond}{code}}。

\begin{texcode}[]{}
  \begin{algorithm}[H]
    \SetAlgoLined
    \KwData{this text}
    \KwResult{how to write algorithm with \LaTeX2e }
  
    initialization\;
    \While{not at end of this document}{
      read current\;
      \eIf{understand}{
        go to next section\;
        current section becomes this one\;
        }{
        go back to the beginning of current section\;
        }
      }
    \caption{How to write algorithms}
  \end{algorithm}
\end{texcode}

\begin{algorithm}[h]
  \caption{Simulation-optimization heuristic}\label{algorithm:SOH}
  \KwData{current period $t$, initial inventory $I_{t-1}$, initial capital $B_{t-1}$, demand samples}
  \KwResult{Optimal order quantity $Q^{\ast}_{t}$}
  $r\leftarrow t$\;
  $\Delta B^{\ast}\leftarrow -\infty$\;
  \While{$\Delta B\leq \Delta B^{\ast}$ and $r\leq T$}{$Q\leftarrow\arg\max_{Q\geq 0}\Delta B^{Q}_{t,r}(I_{t-1},B_{t-1})$\;
  $\Delta B\leftarrow \Delta B^{Q}_{t,r}(I_{t-1},B_{t-1})/(r-t+1)$\;
  \If{$\Delta B\geq \Delta B^{\ast}$}{$Q^{\ast}\leftarrow Q$\;
  $\Delta B^{\ast}\leftarrow \Delta B$\;}
  $r\leftarrow r+1$\;}
\end{algorithm}


\begin{algorithm}[H]
  \KwData{this text}
  \KwResult{how to write algorithm with \LaTeX2e }
  initialization\;
  \While{not at end of this document}{
    read current\;
    \Repeat{this end condition}{
      do these things\;
    }
    \eIf{understand}{
      go to next section\;
      current section becomes this one\;
    }{
      go back to the beginning of current section\;
    }
    \Do{this end condition}{
      do these things\;
    }
  }
  \caption{How to write algorithms}
\end{algorithm}



\xsection{导入代码}{Input codes}

从现有文件导入源码，使用 \clist[columbiablue]{lstlisting} 宏包，本模板预先定义了一基本格式 \clist{sty_basic} 以供使用，可以自定义新的样式。

\begin{texcodeonly}
\lstinputlisting[style=sty_basic,
  language=matlab,
  caption={标准粒子群算法},
  label={listing:a_label}]{./Codes/pso.m}
\end{texcodeonly}
\lstinputlisting[style=sty_basic,
                    language=matlab,
                    caption={标准粒子群算法},
                    label={listing:a_label}]{./Codes/pso.m}
% !TeX root = ../main.tex

\xchapter{参考文献与交叉引用}{Crossreferences}

注意，当引用时出现 [?] 时，可能有有多种原因造成，首先请确定是否使用了 biber 编译模板，再请注意查看辅助文件是否正常生成，此种问题大部分情况是使用不当而非模板设置造成的。

\xsection{基于 Label 的引用}{Ref labels}

% \xsubsection{引用公式}{Ref equations}

% 引用公式使用 \clist{\eqref{label}} 命令，将引用公式的编号（带括号）：

% \begin{texcode}[]{}
%   如式 \eqref{eqn:c1:diff} 所示
% \end{texcode}

% \xsubsection{引用其他内容}{Ref other labels}

基于 \clist{\label{}} 编号的引用，可以通过 \clist[columbiablue]{cleveref} 提供的 \clist{\cref} 实现，它可以通过识别 label 的类型，自动为引用添加前缀。具体的使用方法可以查看该宏包的文档。

应当这样使用 \clist{\label{TYPE:NAME}}，其中 TYPE 指此标签的类型，应当使用模板提供的或自行指定；NAME 为此标签的名称，可以自行定义，不应当重复。模板为所有举例过的环境预设了类型，其中在\cref{section:Environments for math}中预设了 $12$种数学环境，则其 TYPE 与该环境的「环境」列保持一致，其余类型的定义见下：

\begin{texcode}[]{}
  如\cref{equation:c1:diff} 所示 \\
  如\crefrange{equation:c1:exp}{equation:c1:diff} 所示 \\
  如\cref{equation:c1:diff,equation:c1:exp,equation:c1:mdl:cstr_dis} 所示 \\
  如\cref{figure:icon} 所示，如\cref{subfigure:icon} 所示 \\
  如\cref{table:methodcompare} 所示 \\
  如\cref{algorithm:SOH} 所示 \\
  如\cref{listing:a_label} 所示 \\
  如\cref{theorem:pt} 所示
\end{texcode}

模板只对数学公式设置了多种引用时的样式，并未对其他类型设置，如果有使用的需要，请自行查看手册与模板文件模仿修改。

\xsection{参考文献}{References}

\begin{tcolorbox}[colback=red!5!white,colframe=red!75!black]
  \begin{enumerate}[leftmargin=0.5cm]
    \item 参考文献的著录格式应符合国家标准GB/T 7714-2015《文后参考文献著录规则》。参考文献中每条项目应齐全。
    \item 参考文献里面标点符号：英文文献用半角,中文文献用全角。
  \end{enumerate}
\end{tcolorbox}

\begin{tcolorbox}[colback=blue!5!white,colframe=blue!75!black]
  \begin{enumerate}[leftmargin=0.5cm]
    \item 文后著录的参考文献务必实事求是。论文中引用过的文献必须著录，未引用的文献不得出现。应遵循学术道德规范，避免涉嫌抄袭、剽窃等学术不端行为。
    \item 参考文献一般应是作者亲自考察过的对学位论文有参考价值的文献，除特殊情况外，一般不应间接引用。
    \item 参考文献应有权威性，要注意引用最新的文献。
    \item 参考文献的数量：
    博士学位论文，一般应在80篇以上，其中，期刊文献60篇以上，国外文献30篇以上，均以近5年的文献为主。
    硕士学位论文，一般应在30篇以上，其中，期刊文献不少于20篇，国外文献不少于10篇，均以近5年的文献为主。
  \end{enumerate}
\end{tcolorbox}

模板使用 biblatex 编译参考文献，默认采用顺序编码制，有多种引用参考文献的方式，请自行查阅 biblatex-gbt-7715-2015 宏包，即 \clist{texdoc biblatex-gbt-7714-2015}。

在导言区使用\clist{\addbibresource{path/to/.bib file}}添加参考文献数据库，默认参考文献数据库位于 ./References/ 下。

这里主要举例三种引用方式，使用 \clist{\cite{ref}} 的引用结果作为上标出现在正文中，使用 \clist{\parencite{ref}} 的引用结果作为正文内容（一个名词）出现在正文中，使用 \clist{\footfullcite{ref}} 的引用结果作为脚注内容在脚注部分出现。



\begin{texcode}[]{}
  这是一个比较常见的问题 \cite{barella_situ_2021}。可以发现 xxxx \cite{atta_enhanced_2021}。\\
  文献\parencite{张燕2013电气自动化在电气工程中的应用探讨} 提出了 xxx，文献 \parencite{黄雪芳2012探讨电气工程中自动化技术的应用} 则提出了 yyy。
\end{texcode}

\footnote{这是脚注1}根据 zzz 发布的报告 \footfullcite{王磊2011电气自动化控制设备可靠性探究}，可以发现\footnote{这是脚注2}
% !TeX root = ../main.tex

\xchapter{非正文部分的要求}{Requirements of other parts}

\xsection{摘要}{Abstracts}

\begin{tcolorbox}[colback=blue!5!white,colframe=blue!75!black,title=摘要内容要求]
  内容一般包括：从事这项研究工作的目的和意义；完成的工作（作者独立进行的研究工作及相应结果的概括性叙述）；获得的主要结论（这是摘要的中心内容）。博士学位论文摘要应突出论文的创新点，硕士学位论文摘要应突出论文的新见解。
\end{tcolorbox}

\begin{tcolorbox}[colback=red!5!white,colframe=red!75!black,title=中文摘要主体格式]
  \begin{enumerate}[leftmargin=0.5cm]
    \item 论文摘要由摘要正文、关键词、论文类型、资助申明等部分组成。
    \item 博士学位论文摘要正文为1000字(word)左右，硕士学位论文摘要正文为600字(word)左右。
    \item 摘要中一般不用图、表、化学结构式、非公知公用的符号和术语。
    \item 如果论文的主体工作得到了有关基金资助，应在摘要第一页的页脚处标注：本研究得到某某基金（编号：）资助。（五号）
  \end{enumerate}
\end{tcolorbox}

\begin{tcolorbox}[colback=red!5!white,colframe=red!75!black,title=英文摘要主体格式]
  \begin{enumerate}[leftmargin=0.5cm]
    \item 用词准确，符合语法；
    \item 关键词按相应专业的标准术语写出，尽量从《英语主题词表》中摘选；
    \item 如果论文的主体工作得到了有关基金资助，应用英文在摘要第一页的页脚处标注：本研究得到某某基金（编号：）资助；
    \item 中文摘要和英文摘要均不要求学位申请人及其指导教师签字。
    \item 摘要正文每段开头不空格，每段之间空一行；
  \end{enumerate}
\end{tcolorbox}

\begin{tcolorbox}[colback=blue!5!white,colframe=blue!75!black,title=关键词要求]
  关键词由3～5个组成。关键词应从《汉语主题词表》中摘选，当《汉语主题词表》的词不足以反映主题时，可由申请人设计关键词，但须加注。每一关键词之间用分号分开，最后一个关键词后不打标点符号。由申请人设计的关键词，须在该关键词的右上角标注*，并在该页的页脚处注明“*表示非汉语主题词”。
\end{tcolorbox}

\begin{tcolorbox}[colback=blue!5!white,colframe=blue!75!black,title=论文类型要求]
  论文类型包括：a.理论研究（Theoretical Research）；b.应用基础(Application Fundamentals)；c.应用研究(Application Research)；d.研究报告(Research Report)；e.设计报告(Design Report)；f.案例分析(Case Study)；g.调研报告(Investigation Report)；h.产品研发(Product Development)；i.工程设计(Engineering Design)；j.工程/项目管理(Engineering/Project Management)；k.其它（Others）。
\end{tcolorbox}

\xsection{主要符号表}{Glossary}

\begin{tcolorbox}[colback=red!5!white,colframe=red!75!black]
  \begin{enumerate}[leftmargin=0.5cm]
    \item 如果论文中使用了大量的物理量符号、标志、缩略词、专门计量单位、自定义名词和术语等，应将全文中常用的这些符号及意义列出。如果上述符号和缩略词使用数量不多，可以不设专门的主要符号表，但在论文中出现时须加以说明。
    \item 论文中主要符号应全部采用法定单位，特别要严格执行GB3100～3102—93有关“量和单位”的规定。单位名称的书写，可以采用国际通用符号，也可以用中文名称，但全文应统一，不得两种混用。
    \item 缩略词应列出中英文全称。
    \item 主要符号表正文统一左缩进一个字符。
  \end{enumerate}
\end{tcolorbox}


\xsection{攻读学位期间取得的研究成果}{Achievements}

\begin{tcolorbox}[colback=red!5!white,colframe=red!75!black]
  \begin{enumerate}[leftmargin=0.5cm]
    \item 已发表或已录用的学术论文、已出版的专著/译著、已获授权的专利按参考文献格式列出。
    \item 科研获奖，列出格式为：获奖人（排名情况）．项目名称．奖项名称及等级，发奖机构，获奖时间．。
    \item 与学位论文相关的其它成果参照参考文献格式列出。
    \item 全部研究成果连续编号编排。
    \item 用于双盲评审的论文，只列出已发表的学术论文的篇名、发表刊物名称，必须隐去各类论文检索号、期号、卷号、页码；专利号；日期等。
  \end{enumerate}
\end{tcolorbox}



\xsection{答辩委员会会议决议}{Decisions}

\begin{tcolorbox}[colback=red!5!white,colframe=red!75!black]
  \begin{enumerate}[leftmargin=0.5cm]
    \item 填写内容应与学位（毕业）审批材料中答辩委员会决议书一致。
    \item 无需签名。
    \item 盲审论文仅保留“答辩委员会会议决议” 标题
  \end{enumerate}
\end{tcolorbox}
% !TeX root = ../main.tex

\xchapter{本模板已载入的 Packages}{Loaded Packages}

\begin{table}[h]
  \begin{tabular}{|c|c|c|c|c|c|c|c|}
  \hline
  expl3     & ifxetex    & xcolor   & tikz        & tikz-3dplot & zhnumber    & datetime2  & indentfirst \\ \hline
  setspace  & etoolbox   & xpatch   & xparse      & calc        & ulem        & ifthen     & realboxes   \\ \hline
  blindtext & zhlipsum   & amsmath  & unicode-math     & amsthm      & thmtools    & glossaries & upgreek     \\ \hline
  pifont    & array      & commath  & siunitx     & mathtools   & -    & geometry   & ifoddpage   \\ \hline
  emptypage & pdfpages   & xeCJK    & fontenc     & lmodern     & anyfontsize & mathrsfs   & amsfonts    \\ \hline
  ctex      & ifplatform & graphicx & algorithm2e & tcolorbox   & listings    & enumitem   & footmisc    \\ \hline
  fancyvrb  & fancyhdr   & titlesec & float       & newfloat    & multirow    & longtable  & tabu        \\ \hline
  tabularx  & xltabular  & booktabs & caption     & subcaption  & xurl        & hyperref   & bookmark    \\ \hline
  titletoc  & tocloft    & biblatex & appendix    & cleveref    &             &            &             \\ \hline
  \end{tabular}
\end{table}
\thesisbodyend

% 致谢 修改 Main_Miscellaneous/acknowledegment.tex 中的内容
% Acknowledgement, Rewrite your content in Main_Miscellaneous/acknowledegment.tex
\thesisacknowledegment

% [自动生成] 参考文献 默认使用 References/reference.bib，手动指定请在导言区 \addbibresource 处指定
% [Auto Generate] Bibliography, Default file is References/reference.bib, change argument in \addbibresource for manual specification
\thesisbibliography

% 附录（有几个附录就导入几个文件（不加.tex后缀）），
% Appendi(x/ies), argument should not have the .tex suffix
\thesisappendix{Main_Miscellaneous/appendix_a,
                Main_Miscellaneous/appendix_b}

% 攻读学位期间取得的研究成果
% 添加 [auto] 参数则读取通过 \addachivementresource 添加的数据库，否则读取 Main_Miscellaneous/achievement.tex 的内容 注意盲审时需要手动修改格式
% 自动读取的数据库中 若条目含有 AUTHOR+an = {X=highlight} 则第 X 位作者会被加粗
% Achievement, argument [auto] will load data from database added by \addachivementresource, or the data from Main_Miscellaneous/achievement.tex
% \thesisachivements
\thesisachivements[auto]

% 答辩委员会决议 修改 Main_Miscellaneous/decision.tex 中的内容
% Decision, Rewrite your content in Main_Miscellaneous/decision.tex
\thesisdecision

% [自动生成] 常规评阅人名单 需要手动指定两个数字作为'本学位论文共接受 {#1} 位专家评阅，其中常规评阅人 {2}名'内容参数
% [Auto Generate] Reviewer List, set two number as the content in this page 
\thesisreviewers{7}{5}

% [自动生成] 独创性声明
% [Auto Generate] Originality Declaration
\thesisdeclarations

\end{document}